% ASSERT_CONVENTION: natural_units=natural, metric_signature=mostly_minus, fourier_convention=physics, coupling_convention=alpha_s, generator_normalization=T(fund)=1/2, covariant_derivative_sign=D_mu=partial_mu-igA
%
% Lecture 2: Minimal N=1 SUSY Introduction
% Part of: The Dualised Standard Model -- Part III Lecture Notes
%
% SUSY-specific conventions (Wess-Bagger):
%   Weyl spinors: dotted/undotted, alpha = 1,2
%   Superspace metric: epsilon^{12} = +1, epsilon_{12} = -1
%   Superpotential dimension: [W] = [mass]^3
%   R-charge: R[Q] = (N_f - N_c)/N_f, R[W] = 2
%   N_f counting: Dirac flavour pairs (Q_i, Q-tilde_i)
%
\documentclass[11pt,a4paper]{article}
% ASSERT_CONVENTION: natural_units=natural, metric_signature=mostly_minus, fourier_convention=physics, coupling_convention=alpha_s, generator_normalization=T(fund)=1/2, covariant_derivative_sign=D_mu=partial_mu-igA
%
% CONVENTIONS (Seiberg hep-th/9411149)
% Metric: (+,-,-,-) (mostly minus)
% T(fund) = 1/2 per Weyl fermion
% b_0 = (11N_c - 2N_f)/3 for N_f Dirac flavours
% D_mu = partial_mu - ig A_mu^a T^a
% A(fund) = 1 (cubic anomaly coefficient per fundamental Weyl fermion)
% alpha_s = g_s^2 / (4 pi)
% R[Q] = (N_f - N_c)/N_f; R[W] = 2
% N_f counts Dirac flavour pairs (each = 2 Weyl fermions)
%
% This file is \input{} by all lecture files.

%% ---------- Document class ----------
% (loaded by the wrapper; preamble only provides packages and macros)

%% ---------- Standard packages ----------
\usepackage{amsmath,amssymb,amsthm,mathtools}
\usepackage{hyperref}
\usepackage[capitalise,noabbrev]{cleveref}
\usepackage{enumitem}
\usepackage{booktabs}
\usepackage{array}
\usepackage{xcolor}
\usepackage{tikz}
\usetikzlibrary{arrows.meta,decorations.markings,positioning,calc}
\usepackage{fancybox}
\usepackage{tcolorbox}
\tcbuselibrary{skins,breakable}
\usepackage{slashed}              % Feynman slash notation
\usepackage{cite}                 % compressed citation lists

%% ---------- Hyperref configuration ----------
\hypersetup{
  colorlinks=true,
  linkcolor=blue!70!black,
  citecolor=green!50!black,
  urlcolor=blue!80!black,
  pdfauthor={Part III Lecture Notes},
  pdftitle={The Dualised Standard Model}
}

%% ---------- Theorem environments ----------
\theoremstyle{plain}
\newtheorem{theorem}{Theorem}[section]
\newtheorem{lemma}[theorem]{Lemma}
\newtheorem{proposition}[theorem]{Proposition}
\newtheorem{corollary}[theorem]{Corollary}

\theoremstyle{definition}
\newtheorem{definition}[theorem]{Definition}
\newtheorem{example}[theorem]{Example}
\newtheorem{exercise}{Exercise}[section]

\theoremstyle{remark}
\newtheorem{remark}[theorem]{Remark}

%% ---------- Physics macros: gauge groups and representations ----------
\newcommand{\Nc}{N_{\!c}}
\newcommand{\Nf}{N_{\!f}}
\newcommand{\SU}[1]{\mathrm{SU}(#1)}
\newcommand{\UU}[1]{\mathrm{U}(#1)}

% Representations
\newcommand{\fund}{\square}
\newcommand{\antifund}{\overline{\square}}
\newcommand{\adj}{\mathrm{adj}}
\newcommand{\antisym}{\wedge^{\!2}}        % antisymmetric rank-2
\newcommand{\sym}{\mathrm{S}_2}           % symmetric rank-2

%% ---------- Physics macros: group theory constants ----------
% Dynkin index: Tr(T^a T^b) = T(R) delta^{ab}
\newcommand{\Tfund}{\tfrac{1}{2}}          % T(fund) = 1/2 per Weyl fermion
\newcommand{\Tadj}{\Nc}                    % T(adj) = N_c
\newcommand{\Cfund}{\frac{\Nc^2 - 1}{2\Nc}}  % C_2(fund)
\newcommand{\Cadj}{\Nc}                    % C_2(adj)

% Cubic anomaly coefficient: A(R) = Tr_R[T^a {T^b, T^c}] with A(fund) = 1
\newcommand{\Afund}{1}

%% ---------- Physics macros: beta function ----------
\newcommand{\bzero}{b_0}
\newcommand{\bone}{b_1}
\newcommand{\alphas}{\alpha_s}
\newcommand{\gs}{g_s}

%% ---------- Physics macros: SUSY / R-charge ----------
\newcommand{\Rcharge}[1]{R[#1]}
\newcommand{\Wsup}{W}                      % superpotential

%% ---------- Physics macros: covariant derivative and field strength ----------
\newcommand{\Dmu}{D_\mu}
\newcommand{\Fmn}{F_{\mu\nu}}
\newcommand{\Ftilde}{\widetilde{F}_{\mu\nu}}

%% ---------- Physics macros: common abbreviations ----------
\newcommand{\ie}{\textit{i.e.}}
\newcommand{\eg}{\textit{e.g.}}
\newcommand{\cf}{\textit{cf.}}
\newcommand{\IR}{\ensuremath{\mathrm{IR}}}
\newcommand{\UV}{\ensuremath{\mathrm{UV}}}
\newcommand{\AF}{\ensuremath{\mathrm{AF}}}
\newcommand{\BZ}{\ensuremath{\mathrm{BZ}}}
\newcommand{\NSVZ}{\ensuremath{\mathrm{NSVZ}}}
\newcommand{\SQCD}{\ensuremath{\mathrm{SQCD}}}
\newcommand{\QCD}{\ensuremath{\mathrm{QCD}}}
\newcommand{\SM}{\ensuremath{\mathrm{SM}}}
\newcommand{\Tr}{\mathrm{Tr}}

%% ---------- Convention box macro ----------
\newtcolorbox{conventionbox}[1][]{%
  enhanced,
  colback=blue!5!white,
  colframe=blue!60!black,
  fonttitle=\bfseries,
  title={Conventions},
  #1
}

%% ---------- Comparison table style ----------
\newtcolorbox{comparisontable}[1][]{%
  enhanced,
  colback=orange!5!white,
  colframe=orange!70!black,
  fonttitle=\bfseries,
  title={Comparison},
  #1
}

%% ---------- Warning / important box ----------
\newtcolorbox{warningbox}[1][]{%
  enhanced,
  colback=red!5!white,
  colframe=red!70!black,
  fonttitle=\bfseries,
  title={Important},
  #1
}

%% ---------- Preview / forward-reference box ----------
\newtcolorbox{previewbox}[1][]{%
  enhanced,
  colback=green!5!white,
  colframe=green!50!black,
  fonttitle=\bfseries,
  title={Preview},
  #1
}

%% ---------- Page layout ----------
\usepackage[margin=2.5cm]{geometry}
\setlength{\parskip}{0.4em}
\setlength{\parindent}{0em}

%% ---------- Section formatting ----------
\usepackage{titlesec}
\titleformat{\section}{\Large\bfseries}{\thesection}{1em}{}
\titleformat{\subsection}{\large\bfseries}{\thesubsection}{1em}{}

%% ---------- End of preamble ----------


\title{\textbf{Lecture 2: Minimal $\mathcal{N}=1$ Supersymmetry}}
\author{The Dualised Standard Model --- Part III Lecture Notes}
\date{Lent Term 2027}

\begin{document}
\maketitle

\tableofcontents
\newpage

%% ========================================================================
%%  CONVENTIONS
%% ========================================================================

\begin{conventionbox}[title={Conventions for Lecture 2}]
All conventions from Lecture~1 remain in force.  Additional conventions
for supersymmetry (following Wess--Bagger \cite{WessBagger} and
Seiberg~\cite{Seiberg1995}):

\renewcommand{\arraystretch}{1.3}
\begin{tabular}{@{}l l l@{}}
\toprule
\textbf{Quantity} & \textbf{Convention} & \textbf{Comment} \\
\midrule
Units & $\hbar = c = 1$ (natural) & \\
Metric & $\eta_{\mu\nu} = \mathrm{diag}(+1,-1,-1,-1)$ & mostly minus \\
Dynkin index & $\Tr(T^a T^b) = \Tfund\, \delta^{ab}$ & fundamental of $\SU{\Nc}$ \\
Covariant derivative & $\Dmu = \partial_\mu - i\gs A_\mu^a T^a$ & \\
Weyl spinors & $\alpha = 1,2$ (undotted); $\dot\alpha = 1,2$ (dotted) & two-component \\
Superspace metric & $\epsilon^{12} = +1$, $\epsilon_{12} = -1$ & raises/lowers spinor indices \\
$\sigma$-matrices & $\sigma^\mu = (\mathbf{1}, \vec{\sigma})$,
  $\bar\sigma^\mu = (\mathbf{1}, -\vec{\sigma})$ & $\sigma^0 = \bar\sigma^0 = \mathbf{1}_{2\times 2}$ \\
Grassmann coordinates & $\theta^\alpha$, $\bar\theta_{\dot\alpha}$ & $[\theta] = [\text{mass}]^{-1/2}$ \\
Superpotential & $[\Wsup] = [\text{mass}]^3$ & holomorphic in chiral superfields \\
$\Nf$ counting & $\Nf$ = Dirac pairs & each pair $= (Q_i, \widetilde{Q}_i) = 2$ Weyl \\
$R$-charge & $\Rcharge{Q} = (\Nf - \Nc)/\Nf$ & derived from anomaly cancellation \\
\bottomrule
\end{tabular}
\end{conventionbox}


%% ========================================================================
%%  SECTION 1: Motivation
%% ========================================================================
\section{Why supersymmetry?}
\label{sec:motivation}

Supersymmetry (SUSY) has not been observed in nature---no superpartner of
any known particle has been found.  Why, then, devote a lecture to it?
The answer is \emph{not} phenomenological but \emph{theoretical}: SUSY
provides a set of exact analytical tools that are unavailable in
non-supersymmetric gauge theories.  In particular:
\begin{enumerate}[label=(\roman*)]
  \item \textbf{Holomorphy and non-renormalisation.}  The superpotential,
    which controls the interactions of a SUSY theory, is a holomorphic
    function of the fields.  It receives \emph{no} perturbative quantum
    corrections to \emph{all orders} in perturbation theory.  This
    drastically constrains the dynamics.
  \item \textbf{Exact beta functions.}  The NSVZ exact beta function
    relates the running of the gauge coupling to the anomalous dimensions
    of matter fields.  At a conformal fixed point, this gives exact
    results for scaling dimensions.
  \item \textbf{Electric--magnetic duality.}  SUSY is the \emph{only}
    known framework in which electric--magnetic duality between strongly-
    and weakly-coupled gauge theories can be \emph{proved}---that is,
    verified by exact anomaly matching, moduli space matching, and
    consistency under deformations.  This is Seiberg duality
    \cite{Seiberg1995}, the centrepiece of Lecture~3.
\end{enumerate}

\begin{remark}[Strategy for the course]
  The logic of this course is: (1)~introduce the minimal SUSY formalism
  (this lecture); (2)~\emph{prove} Seiberg duality in the supersymmetric
  setting (Lecture~3); (3)~understand its microscopic origin from $\mathcal{N}=2$
  theory (Lecture~4); (4)~extend the duality, with appropriate caveats,
  to the non-supersymmetric Standard Model (Lectures~5--8).
  Understanding SUSY duality is a \emph{prerequisite} for the conjectural
  non-SUSY extension that is the main goal of the course.
\end{remark}


%% ========================================================================
%%  SECTION 2: The SUSY Algebra
%% ========================================================================
\section{The $\mathcal{N}=1$ supersymmetry algebra}
\label{sec:susy-algebra}

\begin{definition}[Supersymmetry]
\label{def:susy}
A \textbf{supersymmetry} is a symmetry whose generators $Q_\alpha$
carry spin~$1/2$.  Unlike the bosonic generators of the Poincar\'e group
(translations $P_\mu$, Lorentz transformations $M_{\mu\nu}$), the
supercharges $Q_\alpha$ are \emph{fermionic}: they anticommute rather
than commute, and they transform bosons into fermions and vice versa.
\end{definition}

In four spacetime dimensions, the minimal ($\mathcal{N}=1$) supercharges
form a two-component Weyl spinor $Q_\alpha$ (with $\alpha = 1,2$) and
its conjugate $\bar{Q}_{\dot\alpha}$ (with $\dot\alpha = 1,2$).  Here
$\alpha$ is an \emph{undotted} spinor index transforming under
$\SU{2}_L$, and $\dot\alpha$ is a \emph{dotted} index transforming under
$\SU{2}_R$, where $\SU{2}_L \times \SU{2}_R$ is the double cover of the
Lorentz group $\mathrm{SO}(3,1)$.

\begin{definition}[Sigma matrices]
\label{def:sigma}
The \textbf{sigma matrices} $\sigma^\mu_{\alpha\dot\alpha}$ and
$\bar\sigma^{\mu\,\dot\alpha\alpha}$ connect spinor and vector
indices.  Explicitly:
\begin{equation}
  \sigma^\mu = (\mathbf{1}, \sigma^1, \sigma^2, \sigma^3), \qquad
  \bar\sigma^\mu = (\mathbf{1}, -\sigma^1, -\sigma^2, -\sigma^3),
  \label{eq:sigma-def}
\end{equation}
where $\sigma^i$ are the standard Pauli matrices.  These satisfy
$\sigma^\mu \bar\sigma^\nu + \sigma^\nu \bar\sigma^\mu
= 2\eta^{\mu\nu}\mathbf{1}$.
\end{definition}

The $\mathcal{N}=1$ SUSY algebra extends the Poincar\'e algebra by the
anticommutation relations:
\begin{equation}
  \boxed{
    \{Q_\alpha,\, \bar{Q}_{\dot\alpha}\}
    = 2\,\sigma^\mu_{\alpha\dot\alpha}\, P_\mu\,,
  }
  \label{eq:susy-algebra}
\end{equation}
together with
\begin{equation}
  \{Q_\alpha, Q_\beta\} = 0, \qquad
  \{\bar{Q}_{\dot\alpha}, \bar{Q}_{\dot\beta}\} = 0, \qquad
  [Q_\alpha, P_\mu] = 0.
  \label{eq:susy-algebra-rest}
\end{equation}

\begin{remark}[Key consequences]
\label{rem:susy-consequences}
Three immediate consequences of eq.~\eqref{eq:susy-algebra}:
\begin{enumerate}[label=(\alph*)]
  \item Since $Q_\alpha$ carries spin~$1/2$, each SUSY multiplet
    contains particles differing by $\Delta s = 1/2$.  SUSY relates
    bosons and fermions within the same multiplet.
  \item Every SUSY multiplet has \textbf{equal numbers of bosonic and
    fermionic degrees of freedom}.  This follows from
    $\Tr[(-1)^F] = 0$ in any representation of the algebra,
    where $F$ is the fermion number.
  \item The Hamiltonian is $H = P^0 = \tfrac{1}{4}(Q_1 \bar{Q}_{\dot 1}
    + \bar{Q}_{\dot 1} Q_1 + Q_2 \bar{Q}_{\dot 2}
    + \bar{Q}_{\dot 2} Q_2)$, which is a sum of squares.  Hence
    $\langle H \rangle \geq 0$ in any state, and a supersymmetric
    vacuum has $\langle H \rangle = 0$.
\end{enumerate}
\end{remark}


%% ========================================================================
%%  SECTION 3: Chiral Superfields
%% ========================================================================
\section{Chiral superfields}
\label{sec:chiral}

To construct SUSY-invariant Lagrangians efficiently, we use the
\emph{superfield} formalism, which packages the components of a SUSY
multiplet into a single object living in \emph{superspace}.

\begin{definition}[Superspace]
\label{def:superspace}
\textbf{Superspace} extends ordinary spacetime $x^\mu$ by adjoining
anticommuting (Grassmann) coordinates:
\begin{equation}
  (x^\mu,\; \theta^\alpha,\; \bar\theta_{\dot\alpha}),
  \label{eq:superspace-coords}
\end{equation}
where $\theta^\alpha$ and $\bar\theta_{\dot\alpha}$ are two-component
Grassmann spinors satisfying $\{\theta^\alpha, \theta^\beta\} = 0$ and
$\{\bar\theta_{\dot\alpha}, \bar\theta_{\dot\beta}\} = 0$.  Since each
$\theta^\alpha$ squares to zero, any function of $\theta$ has a
\emph{finite} Taylor expansion terminating at $\theta^2 \equiv
\theta^\alpha \theta_\alpha$.
\end{definition}

\begin{definition}[Grassmann integration]
\label{def:grassmann-int}
Integration over Grassmann variables is defined by the Berezin rules:
\begin{equation}
  \int d\theta\; 1 = 0, \qquad \int d\theta\; \theta = 1.
  \label{eq:berezin}
\end{equation}
The key consequence is that $\int d^2\theta\; f(\theta)$ \textbf{extracts
the $\theta^2$ component} of $f$.  Similarly,
$\int d^2\bar\theta$ extracts the $\bar\theta^2$ component, and
$\int d^4\theta \equiv \int d^2\theta\, d^2\bar\theta$ extracts the
$\theta^2\bar\theta^2$ component.
\end{definition}

\begin{definition}[Chiral superfield]
\label{def:chiral-superfield}
A \textbf{chiral superfield} $\Phi$ is a superfield satisfying the
constraint $\bar{D}_{\dot\alpha} \Phi = 0$, where $\bar{D}_{\dot\alpha}$
is the supercovariant derivative.  This constraint is solved by writing
$\Phi$ as a function of the \emph{chiral coordinate}
$y^\mu = x^\mu + i\theta\sigma^\mu\bar\theta$:
\begin{equation}
  \boxed{
    \Phi(y, \theta)
    = \phi(y) + \sqrt{2}\,\theta\,\psi(y) + \theta^2\, F(y)\,,
  }
  \label{eq:chiral-expansion}
\end{equation}
where:
\begin{itemize}
  \item $\phi(x)$ is a \textbf{complex scalar field};
  \item $\psi_\alpha(x)$ is a \textbf{Weyl fermion} (a two-component
    spinor);
  \item $F(x)$ is a complex \textbf{auxiliary field} (no kinetic term,
    eliminated by its equation of motion).
\end{itemize}
\end{definition}

\begin{remark}[Degree-of-freedom counting]
\label{rem:chiral-dof}
Off-shell, a chiral multiplet has:
\begin{itemize}
  \item \textbf{Bosonic:} $\phi$ (complex scalar) $= 2$ real DOF, plus
    $F$ (complex auxiliary) $= 2$ real DOF $\to$ \textbf{4 bosonic}.
  \item \textbf{Fermionic:} $\psi_\alpha$ (complex Weyl spinor) $= 4$ real
    DOF off-shell $\to$ \textbf{4 fermionic}.
\end{itemize}
On-shell, the equation of motion sets $F = -(\partial \Wsup / \partial \phi)^*$,
eliminating 2 bosonic DOF and the Dirac equation reduces $\psi$ to 2 DOF,
giving \textbf{2~bosonic $=$ 2~fermionic} on-shell, as required by SUSY.
\end{remark}

\begin{definition}[Mass dimensions in superspace]
\label{def:superspace-dims}
Since $\{Q, \bar{Q}\} \sim P_\mu$ and $[P_\mu] = [\text{mass}]$,
while $Q$ acts on $\theta$ by translation, we have
$[\theta] = [\text{mass}]^{-1/2}$.  From the expansion
\eqref{eq:chiral-expansion}:
\begin{equation}
  [\Phi] = [\phi] = [\text{mass}], \qquad
  [\psi] = [\text{mass}]^{3/2}, \qquad
  [F] = [\text{mass}]^2.
  \label{eq:chiral-dims}
\end{equation}
Also, $[d^2\theta] = [\text{mass}]$ (since $[d\theta] = [\theta]^{-1}
= [\text{mass}]^{1/2}$) and $[d^4\theta] = [\text{mass}]^2$.
\end{definition}


\subsection{The superpotential}
\label{sec:superpotential}

\begin{definition}[Superpotential]
\label{def:superpotential}
The \textbf{superpotential} $\Wsup(\Phi)$ is a \textbf{holomorphic}
function of chiral superfields---that is, it depends on $\Phi_i$ but
\emph{not} on $\Phi_i^\dagger$.  In a renormalisable theory,
\begin{equation}
  \Wsup(\Phi) = \frac{1}{2}\, m_{ij}\, \Phi_i \Phi_j
  + \frac{1}{3}\, y_{ijk}\, \Phi_i \Phi_j \Phi_k\,,
  \label{eq:superpotential-general}
\end{equation}
where $m_{ij}$ and $y_{ijk}$ are symmetric couplings.  Since $[\Phi] =
[\text{mass}]$ and $[d^2\theta] = [\text{mass}]$, the Lagrangian term
$\int d^2\theta\, \Wsup$ has $[\Wsup] = [\text{mass}]^3$ and the
full term has dimension~4 as required.
\end{definition}

The superpotential contribution to the Lagrangian is:
\begin{equation}
  \mathcal{L}_\Wsup
  = \int d^2\theta\; \Wsup(\Phi) + \text{h.c.}
  = F_i\, \frac{\partial \Wsup}{\partial \phi_i}
  - \frac{1}{2}\, \frac{\partial^2 \Wsup}{\partial \phi_i \partial \phi_j}\,
  \psi_i \psi_j + \text{h.c.}
  \label{eq:Lsuperpotential}
\end{equation}
The first term is the coupling of the auxiliary field $F_i$ to the
derivative of $\Wsup$, and the second generates Yukawa interactions
and fermion masses from $\Wsup$.

\begin{definition}[K\"ahler potential]
\label{def:kahler}
The \textbf{K\"ahler potential} $K(\Phi^\dagger, \Phi)$ is a real
function of chiral superfields and their conjugates.  It determines the
kinetic terms for the matter fields via:
\begin{equation}
  \mathcal{L}_K = \int d^4\theta\; K(\Phi^\dagger, \Phi).
  \label{eq:kahler-lagrangian}
\end{equation}
For a renormalisable theory, the canonical choice is
$K = \Phi_i^\dagger \Phi_i$, which yields the standard kinetic terms
$\partial_\mu \phi_i^* \partial^\mu \phi_i + i\bar\psi_i \bar\sigma^\mu
\partial_\mu \psi_i + F_i^* F_i$.
\end{definition}


%% ========================================================================
%%  SECTION 4: Vector Superfields and Gauge Theories
%% ========================================================================
\section{Vector superfields and gauge theories}
\label{sec:vector}

To describe gauge fields in the superspace formalism, we need a second
type of superfield.

\begin{definition}[Vector superfield]
\label{def:vector-superfield}
A \textbf{vector superfield} $V$ is a superfield satisfying the reality
condition $V = V^\dagger$.  Under a gauge transformation with chiral
parameter $\Lambda$, the vector superfield transforms as
$e^{2gV} \to e^{-2ig\Lambda^\dagger} e^{2gV} e^{2ig\Lambda}$.
Most of the components of $V$ can be gauged away.
\end{definition}

\begin{definition}[Wess--Zumino gauge]
\label{def:wz-gauge}
The \textbf{Wess--Zumino (WZ) gauge} is a gauge choice that eliminates
the redundant components of $V$, leaving:
\begin{equation}
  \boxed{
    V_{\mathrm{WZ}} = -\theta \sigma^\mu \bar\theta\, A_\mu
    + i\,\theta^2\, \bar\theta \bar\lambda
    - i\,\bar\theta^2\, \theta \lambda
    + \tfrac{1}{2}\,\theta^2 \bar\theta^2\, D\,,
  }
  \label{eq:vector-WZ}
\end{equation}
where the component fields are:
\begin{itemize}
  \item $A_\mu(x)$: the \textbf{gauge boson} (spin-1);
  \item $\lambda_\alpha(x)$: the \textbf{gaugino}, a Weyl fermion in
    the \textbf{adjoint} representation of the gauge group;
  \item $D(x)$: a real \textbf{auxiliary field} (no kinetic term).
\end{itemize}
\end{definition}

\begin{remark}[Degree-of-freedom counting for vector multiplet]
\label{rem:vector-dof}
Off-shell, a vector multiplet in WZ gauge has:
\begin{itemize}
  \item \textbf{Bosonic:} $A_\mu$ has 3~polarisations (4 components minus
    1 residual gauge DOF in WZ gauge), plus $D$ (1 real auxiliary)
    $\to$ \textbf{4 bosonic}.
  \item \textbf{Fermionic:} $\lambda_\alpha$ (complex Weyl fermion)
    $= 4$ real DOF off-shell $\to$ \textbf{4 fermionic}.
\end{itemize}
On-shell, eliminating $D$ and imposing the equations of motion gives
\textbf{2~bosonic} (transverse polarisations) \textbf{$=$ 2~fermionic}
(on-shell gaugino).  The matching $2_B = 2_F$ per multiplet holds in
both the chiral and vector cases.
\end{remark}


\subsection{The field-strength superfield}
\label{sec:field-strength}

\begin{definition}[Field-strength superfield]
\label{def:Walpha}
The \textbf{field-strength superfield} $W_\alpha$ is defined by
\begin{equation}
  W_\alpha = -\frac{1}{4}\, \bar{D}^2 D_\alpha V\,,
  \label{eq:Walpha-def}
\end{equation}
where $D_\alpha$ and $\bar{D}_{\dot\alpha}$ are the supercovariant
derivatives.  Crucially, $W_\alpha$ is itself a \textbf{chiral} superfield
($\bar{D}_{\dot\beta} W_\alpha = 0$), even though it is built from the
real superfield $V$.  Its lowest components are:
\begin{equation}
  W_\alpha = -i\lambda_\alpha + \theta_\alpha\, D
  - \frac{i}{2}(\sigma^\mu \bar\sigma^\nu)_\alpha{}^\beta\, \theta_\beta\, F_{\mu\nu}
  + \theta^2\, \sigma^\mu_{\alpha\dot\alpha}\, \partial_\mu \bar\lambda^{\dot\alpha}.
  \label{eq:Walpha-components}
\end{equation}
\end{definition}

The gauge kinetic Lagrangian is then:
\begin{equation}
  \mathcal{L}_{\text{gauge}}
  = \frac{1}{4g^2} \int d^2\theta\; \Tr\!\bigl(W^\alpha W_\alpha\bigr) + \text{h.c.}
  = \Tr\!\Bigl[
    -\frac{1}{4g^2}\, F_{\mu\nu} F^{\mu\nu}
    + \frac{i}{g^2}\, \bar\lambda \bar\sigma^\mu D_\mu \lambda
    + \frac{1}{2g^2}\, D^2
  \Bigr].
  \label{eq:L-gauge}
\end{equation}
The trace is over gauge indices.  This contains the standard Yang--Mills
kinetic term, the gaugino kinetic term, and the auxiliary $D$-field.

\subsection{Coupling matter to gauge fields}
\label{sec:gauge-matter}

A chiral superfield $\Phi$ in representation $R$ of the gauge group
couples to the vector superfield via the K\"ahler potential:
\begin{equation}
  K = \Phi^\dagger\, e^{2gV}\, \Phi\,,
  \label{eq:kahler-gauge}
\end{equation}
where $V = V^a T^a_R$ with $T^a_R$ the generators in representation $R$.
Expanding in WZ gauge, this produces the gauge-covariant kinetic terms
for $\phi$ and $\psi$, the gauge interactions, and a coupling of $D$
to the scalar fields:
\begin{equation}
  \int d^4\theta\; \Phi^\dagger e^{2gV} \Phi
  \;\supset\;
  (D_\mu \phi)^* (D^\mu \phi) + i\bar\psi \bar\sigma^\mu D_\mu \psi
  + |F|^2 + g\, \phi^* D^a T^a \phi
  + \sqrt{2}\, g\, (\phi^* \lambda^a T^a \psi + \text{h.c.})\,,
  \label{eq:kahler-gauge-components}
\end{equation}
where $D_\mu = \partial_\mu - ig A_\mu^a T^a$ is the gauge-covariant
derivative with our convention.


%% ========================================================================
%%  SECTION 5: The N=1 Lagrangian and Scalar Potential
%% ========================================================================
\section{The $\mathcal{N}=1$ Lagrangian and scalar potential}
\label{sec:lagrangian}

Combining the ingredients from
\S\S\ref{sec:chiral}--\ref{sec:vector}, the general $\mathcal{N}=1$
gauge theory Lagrangian is:
\begin{equation}
  \boxed{
    \mathcal{L} = \int d^4\theta\;
    \Phi_i^\dagger\, e^{2gV}\, \Phi_i
    + \left[\int d^2\theta\;\Bigl(
      \Wsup(\Phi) + \frac{\tau}{16\pi i}\,
      \Tr(W^\alpha W_\alpha)
    \Bigr) + \text{h.c.}\right],
  }
  \label{eq:N1-lagrangian}
\end{equation}
where $\tau = \theta_{\text{YM}}/(2\pi) + 4\pi i / g^2$ is the
\emph{complexified gauge coupling}.  For a renormalisable theory, the
K\"ahler potential is canonical: $K = \Phi_i^\dagger e^{2gV} \Phi_i$.

\subsection{Eliminating auxiliary fields}
\label{sec:aux-elim}

The auxiliary fields $F_i$ and $D^a$ have no kinetic terms and are
eliminated by their algebraic equations of motion:
\begin{align}
  F_i^* &= -\frac{\partial \Wsup}{\partial \phi_i}\,,
  \label{eq:F-eom} \\
  D^a &= -g\, \sum_i \phi_i^* T^a \phi_i\,.
  \label{eq:D-eom}
\end{align}

Substituting back into the Lagrangian yields the \textbf{scalar
potential}:

\begin{definition}[Scalar potential]
\label{def:scalar-potential}
The $\mathcal{N}=1$ scalar potential is
\begin{equation}
  \boxed{
    V = V_F + V_D\,,
  }
  \label{eq:V-total}
\end{equation}
where the \textbf{$F$-term potential} is
\begin{equation}
  \boxed{
    V_F = \sum_i \left|\frac{\partial \Wsup}{\partial \phi_i}\right|^2,
  }
  \label{eq:VF}
\end{equation}
and the \textbf{$D$-term potential} is
\begin{equation}
  \boxed{
    V_D = \frac{1}{2}\sum_a
    \Bigl(g\, \sum_i \phi_i^* T^a \phi_i\Bigr)^2.
  }
  \label{eq:VD}
\end{equation}
\end{definition}

\begin{warningbox}[title={$V \geq 0$: a special feature of SUSY}]
Both $V_F$ and $V_D$ are manifestly \textbf{non-negative}
($V_F$ is a sum of absolute squares, and $V_D$ is a sum of real squares).
Therefore the total scalar potential satisfies $V \geq 0$.

A \textbf{supersymmetric vacuum} has $V = 0$, which requires
\emph{both} $F_i = 0$ for all $i$ \emph{and} $D^a = 0$ for all $a$.
If such a vacuum exists, it is automatically a global minimum of the
potential---there is no tunnelling to a lower-energy state.  SUSY is
unbroken if and only if $V = 0$ can be achieved.
\end{warningbox}


%% ========================================================================
%%  SECTION 6: Non-Renormalisation Theorem
%% ========================================================================
\section{The non-renormalisation theorem}
\label{sec:nonrenorm}

The most powerful consequence of the holomorphy of the superpotential is:

\begin{theorem}[Non-renormalisation \textnormal{(Seiberg~\cite{Seiberg1993})}]
\label{thm:nonrenorm}
The superpotential $\Wsup(\Phi)$ receives \textbf{no perturbative
quantum corrections}---to all orders in perturbation theory.  That is,
the tree-level superpotential is exact perturbatively.
Non-perturbative corrections (instantons) \emph{can} modify $\Wsup$,
but they are exponentially suppressed as $e^{-8\pi^2/g^2}$.
\end{theorem}

\begin{remark}[Seiberg's spurion argument]
\label{rem:seiberg-argument}
Seiberg's elegant proof \cite{Seiberg1993} proceeds as follows.
Promote all couplings in $\Wsup$ to background chiral superfields.
For example, in $\Wsup = m \Phi^2 + \lambda \Phi^3$, treat $m$ and
$\lambda$ as chiral superfields with specific $R$-charge and global
symmetry assignments.  Then:
\begin{enumerate}[label=(\roman*)]
  \item The quantum-corrected superpotential $\Wsup_{\text{eff}}$ must be a
    \textbf{holomorphic} function of the background superfields
    (by supersymmetry);
  \item it must respect all global symmetries and $R$-symmetry;
  \item it must reduce to $\Wsup_{\text{tree}}$ in the weak-coupling
    limit $\lambda \to 0$.
\end{enumerate}
These three constraints---holomorphy, symmetry, and the weak-coupling
limit---\emph{uniquely fix} $\Wsup_{\text{eff}} = \Wsup_{\text{tree}}$.
No perturbative correction is consistent with all three requirements
simultaneously.  Non-perturbative corrections (involving $e^{-8\pi^2/g^2}$,
which is non-holomorphic in the coupling) can and do arise, but they
vanish in perturbation theory.
\end{remark}

\begin{remark}[What is \emph{not} protected]
The K\"ahler potential $K(\Phi^\dagger, \Phi)$ \emph{is} renormalised
and receives corrections at every loop order.  The kinetic terms and
wave-function renormalisations are \emph{not} protected by holomorphy.
This is why anomalous dimensions $\gamma_i$ are nontrivial even in
SUSY theories.
\end{remark}


\subsection{The NSVZ exact beta function}
\label{sec:NSVZ}

The combination of the non-renormalisation theorem for $\Wsup$ and the
fact that the gauge kinetic function $\tau(g)$ is holomorphic leads to
the \textbf{Novikov--Shifman--Vainshtein--Zakharov (NSVZ)} exact beta
function \cite{NSVZ1983,NSVZ1986}:
\begin{equation}
  \boxed{
    \beta(g) = -\frac{g^3}{16\pi^2}\,
    \frac{3\,T(G) - \sum_i T(R_i)\,(1 - \gamma_i)}
         {1 - T(G)\, g^2/(8\pi^2)}\,,
  }
  \label{eq:NSVZ}
\end{equation}
where $T(G)$ is the Dynkin index of the adjoint representation
($T(G) = \Nc$ for $\SU{\Nc}$), $T(R_i)$ is the Dynkin index
of the matter representation $R_i$, and $\gamma_i$ is the anomalous
dimension of the $i$-th matter field.

\begin{remark}[One-loop verification]
\label{rem:NSVZ-oneloop}
At one loop, $\gamma_i = 0$ and the denominator $\to 1$.
For $\SU{\Nc}$ with $\Nf$ chiral superfield pairs
$(Q_i, \widetilde{Q}_i)$, each in the fundamental (with
$T(\fund) = \Tfund$), the numerator gives:
\begin{equation}
  3\,T(G) - \sum_i T(R_i) = 3\Nc - 2\Nf \times \Tfund = 3\Nc - \Nf\,.
  \label{eq:NSVZ-oneloop-num}
\end{equation}
The factor of $2\Nf$ arises because each Dirac flavour pair contributes
two chiral superfields ($Q$ and $\widetilde{Q}$), each with
$T(\fund) = 1/2$.  Therefore the one-loop SUSY beta function coefficient
is:
\begin{equation}
  b_0^{\SQCD} = 3\Nc - \Nf\,.
  \label{eq:b0-sqcd}
\end{equation}
\textit{Consistency check with Lecture~1:}
For a non-SUSY gauge theory, the one-loop coefficient is
$\bzero = (11\Nc - 2\Nf)/3$.  For \SQCD{}, each gauge boson has
a gaugino partner (contributing $-\Nc/3$ to $\bzero$) and each
quark has a squark (contributing $+\Nf/3$), giving an additional
$-(11\Nc/3 - 2\Nf/3) + (3\Nc - \Nf) = (11\Nc - 2\Nf)/3 - (11\Nc - 2\Nf)/3 + (3\Nc - \Nf)$\ldots.  More directly: the $\mathcal{N}=1$ vector multiplet
contributes $T(G) \times 3 = 3\Nc$ (from gluon + gaugino), and
each chiral multiplet $(Q_i$ or $\widetilde{Q}_i)$ contributes
$-T(\fund) = -1/2$ (from scalar + Weyl fermion).  With $2\Nf$
chiral multiplets: $\bzero^{\SQCD} = 3\Nc - 2\Nf \times (1/2) = 3\Nc - \Nf$.
This matches eq.~\eqref{eq:b0-sqcd} and is consistent with the
non-SUSY formula when the SUSY particle content is substituted.
\end{remark}


%% ========================================================================
%%  SECTION 7: N=1 SQCD -- Field Content and Symmetries
%% ========================================================================
\section{$\mathcal{N}=1$ SQCD: field content and symmetries}
\label{sec:sqcd}

\begin{definition}[$\mathcal{N}=1$ SQCD]
\label{def:sqcd}
$\mathcal{N}=1$ \textbf{supersymmetric QCD (SQCD)} is the $\SU{\Nc}$
gauge theory with $\Nf$ chiral superfields $Q^i$ ($i = 1, \ldots, \Nf$)
in the \textbf{fundamental} representation and $\Nf$ chiral superfields
$\widetilde{Q}_i$ in the \textbf{antifundamental} representation, with
\textbf{zero tree-level superpotential}:
\begin{equation}
  \Wsup_{\text{electric}} = 0\,.
  \label{eq:W-electric}
\end{equation}
Each pair $(Q^i, \widetilde{Q}_i)$ constitutes one Dirac flavour; the
total matter content is $\Nf$ Dirac flavour pairs, consistent with
our counting convention.
\end{definition}

The field content of $\mathcal{N}=1$ SQCD, decomposed into component
fields, is:

\begin{center}
\renewcommand{\arraystretch}{1.3}
\begin{tabular}{@{}l l l l@{}}
\toprule
\textbf{Superfield} & \textbf{Boson} & \textbf{Fermion} & \textbf{$\SU{\Nc}$ rep.} \\
\midrule
$V$ (vector) & $A_\mu$ (gauge boson) & $\lambda$ (gaugino) & adjoint \\
$Q^i$ (chiral) & $\tilde{q}^i$ (squark) & $q^i$ (quark) & fundamental \\
$\widetilde{Q}_i$ (chiral) & $\tilde{\bar{q}}_i$ (anti-squark)
  & $\bar{q}_i$ (anti-quark) & antifundamental \\
\bottomrule
\end{tabular}
\end{center}

\subsection{Global symmetries}
\label{sec:global-sym}

The classical global symmetry of SQCD is:
\begin{equation}
  G_{\text{classical}}
  = \SU{\Nf}_L \times \SU{\Nf}_R \times \UU{1}_B \times \UU{1}_A \times \UU{1}_R\,,
  \label{eq:global-sym-classical}
\end{equation}
where $\SU{\Nf}_L$ rotates the $Q^i$, $\SU{\Nf}_R$ rotates the
$\widetilde{Q}_i$, $\UU{1}_B$ is the baryon number, $\UU{1}_A$ is the
axial symmetry, and $\UU{1}_R$ is the $R$-symmetry (a symmetry that
does not commute with SUSY: it rotates the superspace coordinate
$\theta$).

\begin{definition}[$R$-symmetry]
\label{def:R-symmetry}
An \textbf{$R$-symmetry} is a global $\UU{1}$ symmetry under which the
Grassmann coordinate $\theta$ carries charge $+1$:
$\theta \to e^{i\alpha} \theta$.  By the structure of the superfield
expansion \eqref{eq:chiral-expansion}, the component fields of a
chiral superfield $\Phi$ with $R$-charge $r$ have charges:
\begin{equation}
  \Rcharge{\phi} = r, \qquad
  \Rcharge{\psi} = r - 1, \qquad
  \Rcharge{F} = r - 2.
  \label{eq:R-components}
\end{equation}
The gaugino $\lambda$, being the $\theta$-component of the vector
superfield, has $\Rcharge{\lambda} = +1$.  The superpotential must
carry $\Rcharge{\Wsup} = 2$, since the Lagrangian term
$\int d^2\theta\, \Wsup$ must be $R$-neutral and $\Rcharge{\theta^2} = 2$.
\end{definition}

\begin{warningbox}[title={The axial $\UU{1}_A$ is anomalous}]
Just as in non-supersymmetric QCD (Lecture~1, \S3), the axial symmetry
$\UU{1}_A$ is broken by the ABJ anomaly: the triangle diagram
$\SU{\Nc}^2\,\UU{1}_A$ is nonzero.  The anomaly-free global symmetry is
\begin{equation}
  G_{\text{quantum}}
  = \SU{\Nf}_L \times \SU{\Nf}_R \times \UU{1}_B \times \UU{1}_R\,.
  \label{eq:global-sym-quantum}
\end{equation}
\end{warningbox}


\subsection{Deriving the anomaly-free $R$-charge}
\label{sec:R-charge-derivation}

The $R$-charge assignments of the matter fields are \emph{not} a free
choice.  They are \textbf{uniquely determined} by the requirement that
the $R$-symmetry be non-anomalous---specifically, that the mixed
$\SU{\Nc}^2\, \UU{1}_R$ anomaly vanishes.

\begin{remark}[Why uniquely determined]
\label{rem:R-unique}
In SQCD with $\Wsup = 0$, the only constraint on the $R$-charges is
anomaly cancellation together with the requirement $\Rcharge{Q} =
\Rcharge{\widetilde{Q}}$ (from the $\SU{\Nf}_L \times \SU{\Nf}_R$
symmetry and the absence of a superpotential).  This is a single
equation in a single unknown, so the solution is unique.
\end{remark}

The $\SU{\Nc}^2\,\UU{1}_R$ anomaly coefficient receives contributions
from all fermions charged under $\SU{\Nc}$:
\begin{equation}
  \mathcal{A}\bigl[\SU{\Nc}^2\, \UU{1}_R\bigr]
  = \sum_{\text{fermions}} T(R_f)\, R_f\,,
  \label{eq:anomaly-coefficient}
\end{equation}
where $T(R_f)$ is the Dynkin index of the representation $R_f$
and $R_f$ is the $R$-charge of the fermion.

\textbf{Important:} The $R$-charge entering the anomaly formula is the
\emph{fermion} $R$-charge, not the superfield $R$-charge.
Eq.~\eqref{eq:R-components} gives $R_{\text{fermion}} = R_{\text{superfield}} - 1$
for chiral multiplet fermions.  The gaugino, however, is a component
of the vector multiplet and has $R_{\text{gaugino}} = +1$ directly.

Let $r \equiv \Rcharge{Q} = \Rcharge{\widetilde{Q}}$ denote the
\emph{superfield} $R$-charge of the quarks.  The fermion contributions
to the anomaly are:

\begin{center}
\renewcommand{\arraystretch}{1.3}
\begin{tabular}{@{}l c c c c@{}}
\toprule
\textbf{Fermion} & \textbf{Rep.} & $T(R)$ & $R_{\text{fermion}}$
  & \textbf{Multiplicity} \\
\midrule
Gaugino $\lambda$ & adjoint & $\Nc$ & $+1$ & 1 \\
Quark $\psi_Q$ & fund & $\Tfund$ & $r - 1$ & $\Nf$ \\
Anti-quark $\psi_{\widetilde{Q}}$ & anti-fund & $\Tfund$ & $r - 1$ & $\Nf$ \\
\bottomrule
\end{tabular}
\end{center}

Setting the anomaly to zero:
\begin{align}
  0 &= T(\adj) \times (+1) + \Nf\, T(\fund) \times (r - 1)
       + \Nf\, T(\fund) \times (r - 1) \nonumber \\
    &= \Nc + 2\Nf \times \Tfund \times (r - 1) \nonumber \\
    &= \Nc + \Nf\,(r - 1)\,.
  \label{eq:R-anomaly-cancel}
\end{align}
Solving for $r$:
\begin{equation}
  \boxed{
    \Rcharge{Q} = \Rcharge{\widetilde{Q}}
    = 1 - \frac{\Nc}{\Nf}
    = \frac{\Nf - \Nc}{\Nf}\,.
  }
  \label{eq:R-charge-result}
\end{equation}

\textbf{This is a derived result, not a convention.}  The value
$\Rcharge{Q} = (\Nf - \Nc)/\Nf$ is the unique anomaly-free $R$-charge
assignment for massless SQCD.

\begin{remark}[Verification: $\Rcharge{\Wsup} = 2$]
\label{rem:R-check}
Although the electric theory has $\Wsup = 0$, the magnetic dual (Lecture~3)
has a superpotential
$\Wsup_{\text{mag}} = (1/\mu)\, M\, q\, \widetilde{q}$.  For self-consistency,
$\Rcharge{\Wsup} = 2$ must hold.  The meson field $M^i{}_j \sim Q^i
\widetilde{Q}_j$ has $\Rcharge{M} = 2\Rcharge{Q} = 2(\Nf - \Nc)/\Nf$.
The magnetic quarks have $\Rcharge{q} = \Rcharge{\widetilde{q}} = \Nc/\Nf$
(we will derive this in Lecture~3).  Then:
\begin{equation}
  \Rcharge{M q\widetilde{q}}
  = \frac{2(\Nf - \Nc)}{\Nf} + \frac{2\Nc}{\Nf} = 2 = \Rcharge{\Wsup}.
  \quad\checkmark
  \label{eq:R-check}
\end{equation}
\end{remark}


\subsection{Field content and quantum numbers}
\label{sec:field-content-table}

The full quantum numbers of the electric theory are summarised in the
following table:

\begin{center}
\renewcommand{\arraystretch}{1.3}
\begin{tabular}{@{}l c c c c c@{}}
\toprule
\textbf{Field} & $\SU{\Nc}$ & $\SU{\Nf}_L$ & $\SU{\Nf}_R$ & $\UU{1}_B$ & $\UU{1}_R$ \\
\midrule
$Q$ & $\fund$ & $\fund$ & $\mathbf{1}$ & $+1$ & $\dfrac{\Nf - \Nc}{\Nf}$ \\[8pt]
$\widetilde{Q}$ & $\antifund$ & $\mathbf{1}$ & $\antifund$ & $-1$
  & $\dfrac{\Nf - \Nc}{\Nf}$ \\[8pt]
$\lambda$ (gaugino) & $\adj$ & $\mathbf{1}$ & $\mathbf{1}$ & $0$ & $+1$ \\
\bottomrule
\end{tabular}
\end{center}

\noindent This table, together with the analogous table for the magnetic
theory, forms the basis for the anomaly matching verification of
Seiberg duality in Lecture~3.


%% ========================================================================
%%  SECTION 8: Moduli Space of SQCD
%% ========================================================================
\section{Moduli space of SQCD}
\label{sec:moduli}

\begin{definition}[Moduli space]
\label{def:moduli}
The \textbf{moduli space} of a supersymmetric gauge theory is the space
of inequivalent vacuum configurations satisfying $V = 0$---that is,
the space of field values for which both $F_i = 0$ and $D^a = 0$.
The moduli space is parameterised by the vacuum expectation values of
the scalar fields, modulo gauge equivalence.
\end{definition}

Since $\Wsup = 0$ for SQCD, the $F$-term conditions are automatically
satisfied.  The vacuum structure is determined entirely by the $D$-term
conditions.

\subsection{$D$-flat directions}
\label{sec:D-flat}

The $D$-term conditions for $\SU{\Nc}$ are:
\begin{equation}
  D^a = g\, \bigl(Q^{*\,i}\, T^a\, Q_i
  - \widetilde{Q}_i\, T^a\, \widetilde{Q}^{*\,i}\bigr) = 0
  \qquad \forall\; a = 1, \ldots, \Nc^2 - 1.
  \label{eq:D-flat}
\end{equation}
These conditions require, in matrix notation:
\begin{equation}
  Q^\dagger Q - \widetilde{Q}\, \widetilde{Q}^\dagger
  \propto \mathbf{1}_{\Nc \times \Nc}\,,
  \label{eq:D-flat-matrix}
\end{equation}
where $Q$ is viewed as an $\Nc \times \Nf$ matrix and $\widetilde{Q}$
as an $\Nf \times \Nc$ matrix.  The proportionality to the identity
can be absorbed by $\SU{\Nc}$ gauge transformations.

\subsection{Gauge-invariant coordinates}
\label{sec:gauge-invariant}

The \emph{physical} moduli space is the quotient of the $D$-flat
directions by gauge transformations.  It is most naturally described
in terms of gauge-invariant composite operators:

\begin{definition}[Mesons and baryons]
\label{def:mesons-baryons}
\begin{itemize}
  \item \textbf{Mesons:}
    $M^i{}_j = Q^i \widetilde{Q}_j$,
    an $\Nf \times \Nf$ matrix (colour indices contracted).
  \item \textbf{Baryons} (for $\Nf \geq \Nc$):
    $B^{i_1 \cdots i_{\Nc}}
    = \epsilon^{a_1 \cdots a_{\Nc}} Q^{i_1}_{a_1} \cdots Q^{i_{\Nc}}_{a_{\Nc}}$,
    antisymmetric in the flavour indices.
  \item \textbf{Anti-baryons:}
    $\widetilde{B}_{i_1 \cdots i_{\Nc}}
    = \epsilon_{a_1 \cdots a_{\Nc}} \widetilde{Q}^{a_1}_{i_1}
    \cdots \widetilde{Q}^{a_{\Nc}}_{i_{\Nc}}$.
\end{itemize}
\end{definition}

\subsection{Classical constraints and moduli space dimension}
\label{sec:moduli-dim}

The mesons and baryons are not all independent.  Classically, they
satisfy the constraint:
\begin{equation}
  \det M - B\, \widetilde{B} = 0
  \qquad \text{(classical, for $\Nf \geq \Nc$)}.
  \label{eq:classical-constraint}
\end{equation}
Furthermore, the meson matrix has $\mathrm{rank}(M) \leq \Nc$
classically, since $M^i{}_j = Q^i_a \widetilde{Q}^a_j$ factors through
colour space.

The \textbf{dimension} of the classical moduli space is:
\begin{equation}
  \boxed{
    \dim(\mathcal{M}_{\text{cl}})
    = 2\Nc \Nf - (\Nc^2 - 1)\,,
  }
  \label{eq:moduli-dim}
\end{equation}
which counts the $2\Nc\Nf$ real scalar DOF in $(Q, \widetilde{Q})$
minus the $\Nc^2 - 1$ gauge DOF removed by $\SU{\Nc}$ transformations.

\begin{remark}[Consistency check]
For $\Nf = \Nc = 2$: $\dim = 2 \times 2 \times 2 - 3 = 5$.
The moduli space is parameterised by the $2 \times 2$ meson matrix $M$
(4 complex $= 8$ real entries) minus the classical constraint
$\det M = B \widetilde{B}$, which removes 2 real DOF (one complex
constraint), plus the baryons $B$ and $\widetilde{B}$ (1 complex each
$= 4$ real)---but then the constraint itself must be imposed.  The
counting matches: $(4 + 1 + 1 - 1) \times 2_{\mathbb{R}} - 2_{\text{constraint}} = 8$
real DOF, but the true gauge-invariant dimension is 5 complex $= 10$
real... In fact, the formula \eqref{eq:moduli-dim} gives $2 \times 2
\times 2 - 3 = 5$ \emph{complex} dimensions, or 10 real dimensions.
This is the correct answer: 5 independent complex parameters describe
the vacuum.
\end{remark}

\begin{previewbox}[title={Preview: Quantum Modifications (Lecture 3)}]
The classical constraint \eqref{eq:classical-constraint} is modified by
quantum effects:
\begin{itemize}
  \item For $\Nf = \Nc$: the constraint becomes
    $\det M - B\,\widetilde{B} = \Lambda^{2\Nc}$
    (Seiberg~\cite{Seiberg1995}).
  \item For $\Nf < \Nc$: the Affleck--Dine--Seiberg (ADS) superpotential
    is generated non-perturbatively, lifting the moduli space.
  \item For $\Nf > \Nc$: \textbf{Seiberg duality} provides a dual
    magnetic description.  The moduli space, described by mesons and
    baryons on both sides, matches.  Verifying this matching, together
    with the anomaly matching of all global symmetry anomaly
    coefficients, constitutes the strongest evidence for the duality.
\end{itemize}
\end{previewbox}


%% ========================================================================
%%  SECTION 9: Summary and Preview
%% ========================================================================
\section{Summary and preview of Lecture 3}
\label{sec:summary}

\subsection{Key results}

\begin{enumerate}
  \item The $\mathcal{N}=1$ SUSY algebra extends the Poincar\'e algebra
    by fermionic generators $Q_\alpha$, $\bar{Q}_{\dot\alpha}$ satisfying
    $\{Q_\alpha, \bar{Q}_{\dot\alpha}\} = 2\sigma^\mu_{\alpha\dot\alpha}
    P_\mu$.  Each multiplet has equal bosonic and fermionic DOF.

  \item \textbf{Chiral superfields} $\Phi = \phi + \sqrt{2}\theta\psi +
    \theta^2 F$ describe matter: one complex scalar, one Weyl fermion,
    and one auxiliary field.

  \item \textbf{Vector superfields} in WZ gauge contain a gauge boson
    $A_\mu$, a gaugino $\lambda$, and an auxiliary $D$-field.

  \item The \textbf{superpotential} $\Wsup(\Phi)$ is holomorphic and
    receives no perturbative corrections (non-renormalisation theorem).

  \item The \textbf{scalar potential} $V = V_F + V_D \geq 0$, with
    $V_F = \sum |\partial \Wsup / \partial \phi_i|^2$ and
    $V_D = \frac{1}{2}\sum_a (g\, \phi^* T^a \phi)^2$.
    SUSY vacua have $V = 0$.

  \item The \textbf{NSVZ beta function} gives $b_0^{\SQCD} = 3\Nc - \Nf$
    at one loop.

  \item $\mathcal{N}=1$ \textbf{SQCD} is $\SU{\Nc}$ with $\Nf$ pairs
    $(Q, \widetilde{Q})$ and $\Wsup = 0$.  The anomaly-free
    $R$-charge is $\Rcharge{Q} = (\Nf - \Nc)/\Nf$, \textbf{derived}
    from $\SU{\Nc}^2\,\UU{1}_R$ anomaly cancellation.

  \item The \textbf{moduli space} is parameterised by mesons
    $M^i{}_j = Q^i \widetilde{Q}_j$ and baryons, with classical
    dimension $2\Nc\Nf - (\Nc^2 - 1)$.
\end{enumerate}

\subsection{Preview: Lecture 3}

In Lecture~3, we state Seiberg duality: for $\Nc + 2 \leq \Nf \leq
\tfrac{3}{2}\Nc$, the \IR{} dynamics of $\SU{\Nc}$ SQCD is described by
a \emph{magnetic} $\SU{\Nf - \Nc}$ gauge theory with $\Nf$ dual quarks
$(q, \widetilde{q})$, a gauge-singlet meson $M$, and a superpotential
$\Wsup_{\text{mag}} = (1/\mu)\, M\, q\, \widetilde{q}$.  We will verify
this duality by matching all six independent 't~Hooft anomaly coefficients
and checking consistency under mass deformations, using the tools
developed in Lectures~1 and~2.


%% ========================================================================
%%  BIBLIOGRAPHY
%% ========================================================================
\newpage
\begin{thebibliography}{99}

\bibitem{Seiberg1995}
  N.~Seiberg,
  ``Electric--magnetic duality in supersymmetric non-Abelian gauge theories,''
  Nucl.\ Phys.\ \textbf{B435}, 129 (1995)
  [hep-th/9411149].

\bibitem{Seiberg1993}
  N.~Seiberg,
  ``Naturalness versus supersymmetric non-renormalization theorems,''
  Phys.\ Lett.\ \textbf{B318}, 469 (1993)
  [hep-ph/9309335].

\bibitem{WessBagger}
  J.~Wess and J.~Bagger,
  \textit{Supersymmetry and Supergravity},
  2nd ed.\ (Princeton University Press, 1992).

\bibitem{IntriligatorSeiberg1995}
  K.~Intriligator and N.~Seiberg,
  ``Lectures on Supersymmetric Gauge Theories and Electric--Magnetic
  Duality,''
  Nucl.\ Phys.\ Proc.\ Suppl.\ \textbf{45BC}, 1 (1996)
  [hep-th/9509066].

\bibitem{NSVZ1983}
  V.~A.~Novikov, M.~A.~Shifman, A.~I.~Vainshtein, and V.~I.~Zakharov,
  ``Exact Gell-Mann--Low Function of Supersymmetric Yang--Mills Theories
  from Instanton Calculus,''
  Nucl.\ Phys.\ \textbf{B229}, 381 (1983).

\bibitem{NSVZ1986}
  V.~A.~Novikov, M.~A.~Shifman, A.~I.~Vainshtein, and V.~I.~Zakharov,
  ``Beta Function in Supersymmetric Gauge Theories: Instantons versus
  Traditional Approach,''
  Phys.\ Lett.\ \textbf{B166}, 329 (1986).

\bibitem{Sannino2024}
  G.~Cacciapaglia and F.~Sannino,
  ``The Standard Model as a Dual Gauge Theory,''
  arXiv:2407.17281 [hep-ph] (2024).

\bibitem{PeskinTASI1996}
  M.~E.~Peskin,
  ``Duality in Supersymmetric Yang--Mills Theory,''
  TASI 1996 lectures,
  hep-th/9702094.

\bibitem{Terning2002}
  J.~Terning,
  ``TASI-2002 Lectures: Non-perturbative Supersymmetry,''
  hep-th/0306119 (2003).

\end{thebibliography}

\end{document}
